\section{PDF 的格点计算}\label{sec:numerical}

\subsection{格点矩阵元素}\label{subsec:ME}

在本小节中，我们给出在 CLS $32^3\times96$、$2+1$ 味 clover 费米子系综 H102 上的格点 QCD 计算结果：格距 $a=0.086$~fm，pion 质量 $M_\pi=$ 356~MeV，盒子尺寸 $L\approx 2.7$~fm（$M_\pi L\approx 4.9$）~\cite{Bruno:2014jqa}。价 clover 费米子取 $\kappa_l=0.136865$ 与 $C_{SW}=1.98625$。我们在源/汇涂抹以及在准 PDF 算符 $O_\Gamma$ 中施加两次尺寸为 size=2.5$a$ 的 APE 涂抹~\cite{Hasenfratz:2001hp}，但不作用于费米子传播子。

\begin{table}[htbp]
\begin{center}
\caption{\label{table:p_modes}
NPR 分析中使用的动量模式。括号中的四位数字（以 $2\pi/L$ 为单位）对应三个空间动量分量与能量分量。$z$ 方向可以通过把 Wilson link 设为三个空间方向中的任意一个来选取。因此这些动量模式大致覆盖了 $\mu_R=\sqrt{-(p^R)^2}$ 的三种选择以及若干组 $p^R_z= 2\pi i/L$（$i=0,1,2,...$）。}
\begin{tabular}{c|ccc}
\hline
$a^2\mu_R^2$ & 动量模式 $p^R$ \\
$ [1.109,1.118]$ &(5,2,0,0)   (4,3,1,5)   \\
$ [1.957,1.966]$ &(5,5,0,3)  (6,3,2,4)   (5,4,1,9) \\
$ [2.814,2.832]$ &(6,5,1,10)   (7,4,2,6)   (6,3,0,16)\\
\hline
\end{tabular}
\end{center}
\end{table}

首先，我们将探索 RI/MOM 方案下的非微扰重正化。按照文献~\cite{Chen:2017mzz}，我们使用朗道规范固定的墙源（并把源限制在时间片范围 $t\in[32, 64]$ 内以避免 $t=0$ 处开边界条件带来的边界效应），生成表~\ref{table:p_modes} 所列动量模式的传播子。括号中的四位数字对应以 $2\pi/L$ 为单位的三个空间分量与能量分量。$z$ 方向可通过把 Wilson link 设为三个空间方向中任一方向来选取。于是这些结果近似覆盖了 {$\mu_R=\sqrt{-(p^R)^2}$ 的三个值（2.4、3.2 与 3.9~GeV，对应 $a^2\mu_R^2$=1.1、2.0 与 2.8），以及上限 $p^R_z<\mu_R$ 以内的 $p^R_z=\{0,1,2,...,\}*2\pi /L$}。注意在挑选这些动量模式时，我们要求给定模式的各空间分量彼此不同，并调节 $p_t^R$ 使 $\mu_R$ 保持不变（误差在 2\% 以内）。这些选择使我们能够探究对 $p$ 各分量的依赖；但需注意，由于常规约束 $\frac{\sum_{\mu}a^4p_{\mu}^4}{(\sum_{\mu}a^2p_{\mu}^2)^2}<0.3$ 未被满足，结果可能存在可观的离散化误差。这些离散化误差将在未来研究中考察。

\begin{figure}[htbp]
\includegraphics[width=.48\textwidth]{images/compare_z6.pdf}
\includegraphics[width=.48\textwidth]{images/compare_delta_z6.pdf}
\caption{$z=6a$（$\approx 0.5$~fm）处 NPR 的 $Z=Z_{mp}$（上）与 $\delta Z=Z_{\slashed{p}}-Z_{mp}$（下）随 $p_z^R/\mu_R=1/\sqrt{r}$ 的变化，对应不同的 $\mu_R$。在 $p_z^R=0$ 处，$Z$ 为实数且 $\delta Z/Z$ 小于 5\%。} \label{fig:pz_dependence}
\end{figure}

如式~(\ref{eq:matching_coeff}) 所示，一圈匹配公式主要依赖组合 $r=\mu^2_R/(p_z^R)^2$，而与 $p^R_t$ 无关。图~\ref{fig:pz_dependence} 给出固定 $z\sim 0.5$ fm 处的 $Z_{mp}$ 与 $Z_{\slashed{p}}-Z_{mp}$ 随 $p^R_z/\mu_R=1/\sqrt{r}$ 的变化。由该图可见，在 $p_z^R/\mu_R<0.4$ 范围内，无论实部还是虚部，NPR 因子都只表现出对 $r$ 的依赖，而与 $p_z^R$ 或 $\mu_R$ 的具体取值无关。

\begin{figure}[tbp]
\includegraphics[width=.48\textwidth]{images/Z_pt_dep.pdf}
\caption{最小投影重正化因子的倒数 $1/Z_{mp}$ 及差值 $1/Z_{\slashed{p}}-1/Z_{mp}$ 在相同 $p^R_z$、$\mu_R$ 但不同 $p^R_t$ 下随 $z$ 的变化。叉号对应 $p_t^R=0.9$ GeV，空心方块（圆圈）对应 $p_t^R=2.7$ GeV。多数结果显示对 $p^R_t$ 的依赖较温和。} \label{fig:pt_dependence}
\end{figure}

图~\ref{fig:pt_dependence} 中我们给出了固定 $\mu_R$=3.2 GeV、$p_z^R=1.4$ GeV，以及 $p_t^R$=0.9 GeV 与 2.7 GeV 两个不同取值时，$1/Z_{mp}$ 和 $1/Z_{\slashed{p}}-1/Z_{mp}$ 随 Wilson link 长度 $z$ 的变化。如图所示，两个不同 $p_t^R$ 对应的 $1/Z_{mp}$ 与 $1/Z_{\slashed{p}}-1/Z_{mp}$ 对所有 $z$（同色曲线）都彼此接近。这与一圈匹配公式一致。在 $z<0.3$ fm 处，两种 $p_t^R$ 的 $1/Z_{\slashed{p}}-1/Z_{mp}$ 实部可以略非零，但仍比 $1/Z_{\slashed{p}}$ 小两个数量级。

下面我们取 $p_z^R=0$，并通过改变 $p_z^R$ 来估计来自 $p_z^R$ 依赖的系统不确定度。

在核子矩阵元素的计算中，我们对夸克场采用高斯动量涂抹~\cite{Bali:2016lva}
\begin{align}\label{eq:moms}
\psi(x)\to &S_\text{mom}\psi(x) = \frac{1}{1+6\alpha}\nonumber\\
&\left[\psi(x) + \alpha \sum_j U^{APE}_j(x)e^{ik\hat{e}_j}\psi(x+\hat{e}_j)\right]\, ,
\end{align}
其中 $k$ 是目标动量，{$U^{APE}_j(x)$ 是 $j$ 方向经 APE 涂抹的规范 link}，$\alpha$ 是与传统高斯涂抹一样的可调参数。

这种动量源旨在增大与目标 boost 动量核子的重叠，使我们能够达到比先前工作~\cite{Chen:2017mzz}更高的核子 boost 动量。在这项探索性研究中，虽然我们调节了高斯涂抹半径以便与计算中所用最大动量更好地重叠，但场涂抹在动量空间仍以零动量为中心。改用动量涂抹后，涂抹中心将移到 $O(k)$ 动量处，这将立刻使我们能以更好的信噪比在矩阵元素中达到更高的 boost 动量。本工作中我们使用两个核子 boost 动量 $P_z=n \frac{2\pi}{L}$，$n \in \{4,5\}$，分别对应 1.8 和 2.3~GeV。

\begin{figure*}[htbp]
\includegraphics[width=.49\textwidth,height= 0.35\textwidth]{images/quasipdf-R15-p4.pdf}
\includegraphics[width=.49\textwidth,height= 0.35\textwidth]{images/quasipdf-I15-p4.pdf}\\
\includegraphics[width=.49\textwidth,height= 0.35\textwidth]{images/quasipdf-R15-p5.pdf}
\includegraphics[width=.49\textwidth,height= 0.35\textwidth]{images/quasipdf-I15-p5.pdf}
\caption{非极化 PDF 同位旋矢量核子矩阵元素的实部（左）与虚部（右）在不同动量下随 $z$ 的变化，分别为 $P_z=\frac{8\pi}{L}$=1.8 GeV（上）与 $\frac{10\pi}{L}=$2.3~GeV（下）。我们在裸矩阵元素上施加了 $\{\mu, p_z^R\}=\{3.2, 0\}$ GeV 的 RI/MOM 重正化因子及 $z=0$ 处的归一化，以改善大 $z$ 处的可视性。在给定的正 $z$ 处，数据点略有偏移以显示不同拟合区间得到的基态矩阵元素；从左到右依次为 $t_\text{seq}\in$ [7,9]、[7,10]、[7,11]、[8,11] 与 [9,11]。不同分析在统计误差内彼此一致，而分离度取 7 和 8 的拟合相比其他情形具有较小的不确定度。
}
\label{fig:bareME-tsep}
\end{figure*}

在格点上，我们计算有限 $P_z$ boost 下核子沿 $z$ 方向不随时间变化且非定域的关联函数
\begin{align}
\label{eq:qlat}
\widetilde{h}_\text{lat}(z,P_z,{\Gamma};a^{-1}) =
  \left\langle 0;\vec{P} \right|O_\Gamma(z)
  \left| 0;\vec{P} \right\rangle.
\end{align}
{这里态 $|0; \vec{P}\rangle$ 表示动量为 $\vec{P}=\{0,0,P_z\}$ 的基态（核子）态}。对非极化部分子分布取 $\Gamma=\gamma^t$。

\begin{figure*}[htbp]
\includegraphics[width=.49\textwidth,height= 0.35\textwidth]{images/3ptV2ptZ12tFit7R-jack-comp-p4.pdf}
\includegraphics[width=.49\textwidth,height= 0.35\textwidth]{images/3ptV2ptZ12tFit7I-jack-comp-p4.pdf}\\
\includegraphics[width=.49\textwidth,height= 0.35\textwidth]{images/3ptV2ptZ10tFit7R-jack-comp-p5.pdf}
\includegraphics[width=.49\textwidth,height= 0.35\textwidth]{images/3ptV2ptZ10tFit7I-jack-comp-p5.pdf}
\caption{重正化的非极化 PDF 同位旋矢量核子矩阵元素的实部（左）与虚部（右），${P_z, z}={8\pi/L, 12}$（上）与 ${10\pi/L, 10}$（下），对应 $zP_z\sim9.5$。数据点与采用 $t_\text{seq}\in$ [7, 11] 的拟合所给出的谱带符合得很好。}
\label{fig:bareME-tsep2}
\end{figure*}

\begin{figure}[tbp]
\includegraphics[width=0.48\textwidth,height= 0.3\textwidth]{images/comp-p45-16-Re.pdf}
\includegraphics[width=0.48\textwidth,height= 0.3\textwidth]{images/comp-p45-16-Im.pdf}
\caption{$P_z=$1.8 与 2.3 GeV 的重正化准 PDF 矩阵元素（采用最小投影，$p_z^R=0$、$\mu_R=3.2$ GeV）随 $zP_z$ 的变化。} \label{fig:renormalized_h-tilde}
\end{figure}

随着核子 boost 动量的增大，可以预期激发态贡献会更加严重；因此必须仔细研究激发态污染。为此，我们在五个源—汇分离 $t_\text{seq}\in\{7,8,9,10,11\}\times 0.086$~fm 上计算核子矩阵元素 $\widetilde{h}_{\rm lat}$：$P_z=1.8$ GeV 情形在每批 2005 个规范组态上分别做 $\{4,4,8,8,16\}$ 次测量，$P_z=2.3$ GeV 情形则在 2,000 个组态上测量。我们使用 Chroma 软件包~\cite{Edwards:2004sx}的多格点算法~\cite{Babich:2010qb,Osborn:2010mb}加速夸克传播子的求逆。按照文献~\cite{Bhattacharya:2013ehc}，每个三点（3pt）关联函数 $C_\Gamma^{(3\text{pt})} (t, t_\text{seq})$ 可以分解为（假定源位于 $t= 0$）
\begin{align} \label{eq:fit}
C^\text{3pt}(t,t_\text{seq}; P_z, {\Gamma}) &=
   {\cal A}_0^2 \langle 0 | O_\Gamma | 0 \rangle  e^{-E_0t_\text{seq}} \\
   &+{\cal A}_1^2 \langle 1 | O_\Gamma | 1 \rangle  e^{-E_1t_\text{seq}} \nonumber\\
   &+{\cal A}_1{\cal A}_0 \langle 1 |O_\Gamma | 0 \rangle  e^{-E_1 (t_\text{seq}-t)} e^{-E_0 t} \nonumber\\
   &+{\cal A}_0{\cal A}_1 \langle 0 | O_\Gamma | 1 \rangle  e^{-E_0 (t_\text{seq}-t)} e^{-E_1 t} + \ldots \,,\nonumber
\end{align}
其中 $|n\rangle$（$n > 0$）表示激发态。算符插入在时间 $t$，核子态在汇时刻 $t_\text{seq}$（即源—汇分离）被湮灭。式~(\ref{eq:fit}) 中的谱权重 ${\cal A}_{0,1}$ 与能量 $E_{0,1}$ 可由两点（2pt）关联函数得到：
\begin{align} \label{eq:fit2}
C^\text{2pt}(t_\text{seq}; P_z) &=
   {\cal A}_0^2 e^{-E_0t_\text{seq}} +{\cal A}_1^2 e^{-E_1t_\text{seq}} + \ldots\ .
\end{align}

最终，我们将若干 $t_\text{seq}$ 的 3pt 函数与 2pt 函数做联合拟合，拟合形式取自文献~\cite{Bhattacharya:2013ehc}：
\begin{align}
&\frac{C^\text{3pt}(t,t_\text{seq})}{C^\text{2pt}(t_\text{seq})}= \nn\\
&=\frac{\tilde{h}_{\text{lat}}+C_2(e^{-\Delta Et}+e^{-\Delta E(t_\text{seq}-t)})+C_3e^{-\Delta E t_\text{seq}}}{1+C_1 e^{-\Delta Et_\text{seq}}},\nonumber\\
&C^\text{2pt}(t)=C_0e^{-E_0t}(1+C_1 e^{-\Delta Et}),
\end{align}
其中 $\Delta E= E_1-E_0$。$C_{0,1,2,3}$ 与 $E_{0,1}$ 为自由参数。我们把 3pt/2pt 比值的 $t$ 范围限定为 $t\in[1,t_\text{seq}-1]$，把 2pt 的范围限定为 $t\in[7, 11]$，使拟合的 $\chi^2/d.o.f.$ 为 ${\cal O} (1)$。使用 3pt/2pt 比值而非 3pt 函数本身可以提高拟合的稳定性，特别是当 $t<7$ 的 $C^\text{2pt}(t)$ 也纳入拟合时。

图~\ref{fig:bareME-tsep} 展示了由五个拟合得到的基态核子矩阵元素 $\tilde{h}_{lat}(z,P_z,\gamma_t)$：分别使用分离度 $t_\text{seq}\in$ [7, 9]、[7, 10]、[7,11]、[8,11] 与 [9,11]（数据点对应相同的 $z$，但水平错开以增强可视性）。数据进一步乘以 $\{\mu_R, p_z^R\}=\{3.2, 0\}$ GeV 的重正化因子归一化，并将实部在 $z=0$ 处归一到 1。从图中可以看到，在这些分析中均未出现激发态贡献的明显信号。如果舍弃最小的两个分离度，不确定度会大幅增加。在拟合中，即使 $C_3$ 项在统计上并不显著，我们也保留该项以对其不确定度作出适度估计。

为了比较数据与拟合，我们在图~\ref{fig:bareME-tsep2} 中给出大 $z$ 处 $({P_z, z})=({8\pi/L, 12a})$ 与 $({10\pi/L, 10a})$、$t_\text{seq}\in$ [7,11] 的结果。在这些空间分离处矩阵元素的实部似乎为负。拟合得到的基态贡献以黑色谱带表示。如所见，多数数据可以被拟合很好地描述，因此本文其余部分的结果均采用二态拟合并取区间 $t_\text{seq}\in$ [7, 11]。

两个 $P_z$ 取值的重正化准 PDF 矩阵元素绘于图~\ref{fig:renormalized_h-tilde}，为 $p^R_z=0$、$\mu_R=3.2$ GeV 下随 $zP_z$ 变化的曲线。不同 $P_z$ 的结果在统计不确定度内彼此一致。这表明来自更高扭度效应的幂次修正可能并不可观。

\subsection{系统{不确定度}}

\begin{figure}
\includegraphics[width=0.48\textwidth]{images/error_pzR=0.pdf}
\includegraphics[width=0.48\textwidth]{images/error_pzR=3.pdf}
\caption{\label{fig:error}系统误差的不同来源。详细信息见正文。}
\end{figure}

本小节将考虑四类系统不确定度：傅里叶变换（FT）、非物理标度 $p_z^R$ 与 $\mu_R$、RI/MOM 方案中使用的投影，以及匹配系数的反演。

\begin{figure}[tbp]
\includegraphics[width=.48\textwidth]{images/matching_cutoff.pdf}
\includegraphics[width=.48\textwidth]{images/matching_derivative.pdf}
\caption{\label{fig:matchingPDF} 核子 boost 动量 2.3 GeV 的准 PDF（红色虚线）与最小投影下 $\mu=2$~GeV 的匹配 PDF（黑色实线），RI/MOM 参数为 $p_z^R=0$、$\mu_R=3.2$~GeV。上图与下图分别采用截断法与导数法进行傅里叶变换。匹配策略对 PDF 最终结果有重要影响。}
\end{figure}

下面说明计入这些系统不确定度的细节。

 {1) 傅里叶变换。如图~\ref{fig:renormalized_h-tilde} 所示，$P_z$=2.3 GeV 的 $\widetilde{h}_R(z)$ 在 $z>12a$ 时与零一致。因此在从准 PDF 到 PDF 的标准匹配中，在 $z=12a$ 处截断是合理的。本着这一思路，采用标准 FT 的准 PDF 与匹配 PDF 绘于图~\ref{fig:matchingPDF} 上图，可以看出匹配后的 PDF 表现出振荡行为。文献~\cite{Lin:2017ani}提出了一种“导数”方法来消除这种振荡行为。具体而言，取重正化核子矩阵元素的导数 $\partial_z\widetilde{h}_R(z)$，其傅里叶变换与原矩阵元素的傅里叶变换之间相差一个已知的因子：
%
\begin{equation}
\label{eq:derivative}
\widetilde{q}_R(x) = \int^\infty_{-\infty} \frac{dz}{2\pi} \frac{ie^{i x P_z z}}{x} \partial_z\widetilde{h}_R(z),
\end{equation}
%
前提是当 $|z|\to\infty$ 时 $\widetilde{h}_R(z)\to0$。采用同样的截断，结果示于图~\ref{fig:matchingPDF} 下图，显然振荡行为有所减轻。此外，导数法的结果与标准 FT 法在大部分运动学区域一致，仅在小 $x$ 区不同。这是意料之中的，因为两种方法只在我们做了截断的大 $z$ 区域有差异。我们在图~\ref{fig:error} 中以点划蓝色线展示该差异，同时给出把截断从 $z=10a$ 变到 $14a$ 的误差（绿色点线）。}

\begin{figure}[tbp]
\includegraphics[width=0.45\textwidth]{images/matching_error_pzR=0.pdf}
\includegraphics[width=0.45\textwidth]{images/matching_error_pzR=3.pdf}
\caption{\label{fig:mathcing_error_test} 最小投影下反演匹配公式的效果：黑色实线、红色点线、蓝色点线与绿色点划线分别代表 CT14nnlo PDF、把 CT14nnlo PDF~\cite{Dulat:2015mca} 经逆匹配变为准 PDF、再做匹配回到 PDF、以及迭代匹配所得 PDF 与原始 CT14nnlo PDF 之差。上图（下图）对应 $p_z^R=0$（1.4 GeV）。这些图表明我们所用的反演方法在小 $|x|$ 区可靠性较差。绿色点划线所示的差异已被计入我们的系统误差。}
\end{figure}

 {2) 非物理标度 $p_z^R$ 与 $\mu_R$。RI/MOM 引入了 $p_z^R$ 与 $\mu_R$ 两个非物理标度。原则上，把准 PDF 矩阵元素匹配到光锥 PDF 时，矩阵元素对这两个标度的依赖应与匹配核中的依赖完全相消。然而，由于准 PDF 矩阵元素是在格点上非微扰重正化的，而匹配系数只在微扰论中算到一圈阶，微扰匹配之后会残余这两个标度的依赖。为估计残余的 $p_z^R$ 与 $\mu_R$ 依赖，我们选取 {$p_z^R=0$ GeV 与 $\mu_R=3.2$ GeV} 作为中心值，把 $p_z^R$ 从 $-1.4$ 变到 1.4 GeV（红色虚线），把 $\mu_R$ 从 2.4 变到 3.9 GeV（橙色粗实线）。这些匹配 PDF 之间的差别被视为对非物理标度残余依赖的系统误差，见图~\ref{fig:error}。如图所示，来自 $\mu_R$ 依赖的系统不确定度与其他来源相比很小，但残余的 $p_z^R$ 依赖可能相当可观。}

 {3) 对投影的依赖。本工作讨论了两种投影：最小投影与“$\slashed p$”投影。在 $p_z^R=0$ 时，NPR 因子的投影依赖对所有 $z$ 都小于 5\%，并且在一圈微扰匹配中消失。因此可以预期投影导致的差别同样很小，如图~\ref{fig:error} 中的品红长虚线所示。}

4) 匹配的反演。要从准 PDF 提取 PDF，需要反解因子化公式式~(\ref{eq:fact})。这可以通过改变 $C$ 中 $\alpha_s$ 的符号，再把新的匹配系数与 $\widetilde{q}$ 卷积来实现。更明确地写为
\begin{align}\label{eq:inverse_fact}
q(x,\mu)=&\int_{-\infty}^\infty {dy\over |y|}\: C'\left({x\over y},r,\frac{yP_z}{\mu},\frac{yP_z}{p_z^R}\right)\widetilde{q}(y,P_z, p^R_z,\mu_R)\nonumber\\
&+\mathcal{O}\left({M^2\over P_z^2},{\Lambda_{\text{QCD}}^2\over x^2 P_z^2},\alpha_s^2\right)\,,
\end{align}
其中 $C'=C(\alpha_s\to-\alpha_s)$。我们通过如下操作估计反演因子化公式带来的误差：{从全局分析给出的 PDF 出发~\cite{Dulat:2015mca}，先用式~(\ref{eq:fact}) 再用式~(\ref{eq:inverse_fact}) 返回该 PDF。这一操作理应给出相同的光锥 PDF}。但由于匹配只精确到 $O(\alpha_s)$，两个结果会有差异，该差异很好地估计了来自反演及更高阶修正的系统误差。结果示于图~\ref{fig:mathcing_error_test}，可以发现只有当 $|x|$ 较小时误差才变得可观。{这在预期之内，因为相关的动量标度是 $xP^z$，使得高阶修正在小 $x$ 处变大}。

还有更精细的反演因子化的方法，例如递推程序。然而，如图~\ref{fig:error} 所示，由匹配流程引起的系统误差（青色粗虚线）在大部分区域也小于前两个来源。

如图~\ref{fig:error} 所示，主要的不确定度来自 $p_z^R$ 依赖和 FT。取 $p_z^R=0$ 时，来自不同投影以及反演与匹配的不确定度通常小于 10\%，仅在非常小的 $|x|$ 区域例外。来自 $\mu_R$ 依赖的不确定度则更小。当 $p_z^R=1.4$ GeV 时，这三类来源的不确定度在数值上变大，但仍小于两个主要来源。

